% Options for packages loaded elsewhere
\PassOptionsToPackage{unicode}{hyperref}
\PassOptionsToPackage{hyphens}{url}
\documentclass[
]{article}
\usepackage{xcolor}
\usepackage{amsmath,amssymb}
\setcounter{secnumdepth}{-\maxdimen} % remove section numbering
\usepackage{iftex}
\ifPDFTeX
  \usepackage[T1]{fontenc}
  \usepackage[utf8]{inputenc}
  \usepackage{textcomp} % provide euro and other symbols
\else % if luatex or xetex
  \usepackage{unicode-math} % this also loads fontspec
  \defaultfontfeatures{Scale=MatchLowercase}
  \defaultfontfeatures[\rmfamily]{Ligatures=TeX,Scale=1}
\fi
\usepackage{lmodern}
\ifPDFTeX\else
  % xetex/luatex font selection
\fi
% Use upquote if available, for straight quotes in verbatim environments
\IfFileExists{upquote.sty}{\usepackage{upquote}}{}
\IfFileExists{microtype.sty}{% use microtype if available
  \usepackage[]{microtype}
  \UseMicrotypeSet[protrusion]{basicmath} % disable protrusion for tt fonts
}{}
\makeatletter
\@ifundefined{KOMAClassName}{% if non-KOMA class
  \IfFileExists{parskip.sty}{%
    \usepackage{parskip}
  }{% else
    \setlength{\parindent}{0pt}
    \setlength{\parskip}{6pt plus 2pt minus 1pt}}
}{% if KOMA class
  \KOMAoptions{parskip=half}}
\makeatother
\usepackage{color}
\usepackage{fancyvrb}
\newcommand{\VerbBar}{|}
\newcommand{\VERB}{\Verb[commandchars=\\\{\}]}
\DefineVerbatimEnvironment{Highlighting}{Verbatim}{commandchars=\\\{\}}
% Add ',fontsize=\small' for more characters per line
\newenvironment{Shaded}{}{}
\newcommand{\AlertTok}[1]{\textcolor[rgb]{1.00,0.00,0.00}{\textbf{#1}}}
\newcommand{\AnnotationTok}[1]{\textcolor[rgb]{0.38,0.63,0.69}{\textbf{\textit{#1}}}}
\newcommand{\AttributeTok}[1]{\textcolor[rgb]{0.49,0.56,0.16}{#1}}
\newcommand{\BaseNTok}[1]{\textcolor[rgb]{0.25,0.63,0.44}{#1}}
\newcommand{\BuiltInTok}[1]{\textcolor[rgb]{0.00,0.50,0.00}{#1}}
\newcommand{\CharTok}[1]{\textcolor[rgb]{0.25,0.44,0.63}{#1}}
\newcommand{\CommentTok}[1]{\textcolor[rgb]{0.38,0.63,0.69}{\textit{#1}}}
\newcommand{\CommentVarTok}[1]{\textcolor[rgb]{0.38,0.63,0.69}{\textbf{\textit{#1}}}}
\newcommand{\ConstantTok}[1]{\textcolor[rgb]{0.53,0.00,0.00}{#1}}
\newcommand{\ControlFlowTok}[1]{\textcolor[rgb]{0.00,0.44,0.13}{\textbf{#1}}}
\newcommand{\DataTypeTok}[1]{\textcolor[rgb]{0.56,0.13,0.00}{#1}}
\newcommand{\DecValTok}[1]{\textcolor[rgb]{0.25,0.63,0.44}{#1}}
\newcommand{\DocumentationTok}[1]{\textcolor[rgb]{0.73,0.13,0.13}{\textit{#1}}}
\newcommand{\ErrorTok}[1]{\textcolor[rgb]{1.00,0.00,0.00}{\textbf{#1}}}
\newcommand{\ExtensionTok}[1]{#1}
\newcommand{\FloatTok}[1]{\textcolor[rgb]{0.25,0.63,0.44}{#1}}
\newcommand{\FunctionTok}[1]{\textcolor[rgb]{0.02,0.16,0.49}{#1}}
\newcommand{\ImportTok}[1]{\textcolor[rgb]{0.00,0.50,0.00}{\textbf{#1}}}
\newcommand{\InformationTok}[1]{\textcolor[rgb]{0.38,0.63,0.69}{\textbf{\textit{#1}}}}
\newcommand{\KeywordTok}[1]{\textcolor[rgb]{0.00,0.44,0.13}{\textbf{#1}}}
\newcommand{\NormalTok}[1]{#1}
\newcommand{\OperatorTok}[1]{\textcolor[rgb]{0.40,0.40,0.40}{#1}}
\newcommand{\OtherTok}[1]{\textcolor[rgb]{0.00,0.44,0.13}{#1}}
\newcommand{\PreprocessorTok}[1]{\textcolor[rgb]{0.74,0.48,0.00}{#1}}
\newcommand{\RegionMarkerTok}[1]{#1}
\newcommand{\SpecialCharTok}[1]{\textcolor[rgb]{0.25,0.44,0.63}{#1}}
\newcommand{\SpecialStringTok}[1]{\textcolor[rgb]{0.73,0.40,0.53}{#1}}
\newcommand{\StringTok}[1]{\textcolor[rgb]{0.25,0.44,0.63}{#1}}
\newcommand{\VariableTok}[1]{\textcolor[rgb]{0.10,0.09,0.49}{#1}}
\newcommand{\VerbatimStringTok}[1]{\textcolor[rgb]{0.25,0.44,0.63}{#1}}
\newcommand{\WarningTok}[1]{\textcolor[rgb]{0.38,0.63,0.69}{\textbf{\textit{#1}}}}
\setlength{\emergencystretch}{3em} % prevent overfull lines
\providecommand{\tightlist}{%
  \setlength{\itemsep}{0pt}\setlength{\parskip}{0pt}}
\usepackage{bookmark}
\IfFileExists{xurl.sty}{\usepackage{xurl}}{} % add URL line breaks if available
\urlstyle{same}
\hypersetup{
  hidelinks,
  pdfcreator={LaTeX via pandoc}}

\author{}
\date{}

\begin{document}

\section{Polymarket Insider-Signal Research Brief and
SRS}\label{polymarket-insider-signal-research-brief-and-srs}

Date: March 30, 2026

\subsection{0. Current Project State and Reference Repo
Assessment}\label{0-current-project-state-and-reference-repo-assessment}

\subsubsection{0.1 Current project
state}\label{01-current-project-state}

Your active codebase is a layered Polymarket data collector, not yet a
full detection engine. The live path is \texttt{run\_collector.py} plus
the \texttt{collectors/} modules, with older prototype logic still
present under \texttt{core/} and \texttt{Old-content/}.

Current active architecture:

\begin{Shaded}
\begin{Highlighting}[]
\NormalTok{run\_collector.py}
\NormalTok{  {-}\textgreater{} apply schema}
\NormalTok{  {-}\textgreater{} full Gamma event sync}
\NormalTok{  {-}\textgreater{} full Gamma market sync}
\NormalTok{  {-}\textgreater{} concurrent loops}
\NormalTok{     {-}\textgreater{} websocket market listener}
\NormalTok{     {-}\textgreater{} tier{-}1 order{-}book polling}
\NormalTok{     {-}\textgreater{} tier{-}2 price polling}
\NormalTok{     {-}\textgreater{} trades polling}
\NormalTok{     {-}\textgreater{} metadata resync}
\NormalTok{     {-}\textgreater{} TTL cleanup}
\end{Highlighting}
\end{Shaded}

What the active pipeline already does well:

\begin{itemize}
\tightlist
\item
  Captures market metadata from Gamma at scale.
\item
  Captures a large number of high-frequency snapshots.
\item
  Stores full order book depth for hot markets.
\item
  Archives resolved outcomes for supervised labels.
\item
  Uses a clean collector split by source and task.
\item
  Has enough historical data volume to support serious research.
\end{itemize}

Observed local state from the repo and database:

\begin{itemize}
\tightlist
\item
  SQLite database size is about 11 GB.
\item
  \texttt{snapshots}: 17,605,043 rows.
\item
  \texttt{order\_book\_snapshots}: 1,294,573 rows.
\item
  \texttt{trades}: 31,184 rows.
\item
  \texttt{markets}: 89,478 rows.
\item
  \texttt{events}: 33,383 rows.
\item
  \texttt{market\_resolutions}: 353 rows.
\item
  Snapshot source mix:

  \begin{itemize}
  \tightlist
  \item
    \texttt{ws}: 13,822,437
  \item
    \texttt{gamma}: 2,882,316
  \item
    \texttt{clob}: 900,290
  \end{itemize}
\end{itemize}

Important operational findings:

\begin{itemize}
\tightlist
\item
  The dataset is stale. Latest \texttt{snapshots.captured\_at} is
  \texttt{2026-03-19T03:13:16.609070+00:00}, latest
  \texttt{trades.trade\_time} is
  \texttt{2026-03-19T02:22:32.881000+00:00}, and there are no rows from
  the last hour.
\item
  Logs show the collector was stopped by \texttt{SIGTERM} on March 18,
  2026 during a long market sync, not continuously running.
\item
  The README describes an all-market collector, but
  \texttt{\_load\_tiers()} in \texttt{run\_collector.py} currently
  hard-filters to sports tags only. That means the always-on queues
  exclude most of the event categories relevant to your research vision.
\item
  Among active tier-1 markets, only about 699 of 9,389 are currently
  eligible under the sports filter. Your geopolitics, policy, macro,
  energy, and corporate-information hypotheses are therefore not being
  covered by the real-time path.
\item
  Trade ingestion is materially incomplete for H2. All 31,184 stored
  trades are \texttt{source=\textquotesingle{}ws\textquotesingle{}};
  there are zero \texttt{clob} and zero \texttt{clob\_backfill} rows.
\item
  This likely comes from a schema/API mismatch: the public Polymarket
  \texttt{/trades} endpoint expects \texttt{market} to be a
  comma-separated list of condition IDs, but your collector passes
  \texttt{yes\_token\_id}. Even if the endpoint responds, that is not
  the documented contract.
\item
  Wallet identity data is not being stored anywhere. The current
  \texttt{trades} schema has no \texttt{proxy\_wallet},
  \texttt{transaction\_hash}, \texttt{condition\_id}, \texttt{outcome},
  or profile fields.
\item
  The real-time websocket path subscribes only to
  \texttt{yes\_token\_id} assets, not both yes and no assets. That means
  you are not observing the full two-sided order-flow state of a binary
  market.
\item
  The TTL policy deletes detailed closed-market data after 24 hours.
  That keeps the DB lean, but it is hostile to event studies,
  postmortems, and training labels that need pre-resolution and
  immediate post-resolution context beyond one day.
\item
  \texttt{ml\_pipeline/feature\_builder.py} is still a stub, so the repo
  currently stops at collection/storage rather than modeling, scoring,
  or alerting.
\end{itemize}

Bottom line:

This repo is already a strong market-data foundation, but it is not yet
collecting the identity-aware, event-linked, continuously scored data
required for insider-activity detection or leading-indicator research.

\subsubsection{0.2 What the professor reference repo is
doing}\label{02-what-the-professor-reference-repo-is-doing}

The reference repo at \texttt{Prof-refe/last30days-skill} is not a
market collector. It is a modular multi-source research engine that:

\begin{itemize}
\tightlist
\item
  researches a topic over the last 30 days,
\item
  queries Reddit, X, YouTube, Hacker News, Polymarket, web search, and
  other sources,
\item
  normalizes all sources into a common schema,
\item
  scores them for relevance, recency, and engagement,
\item
  deduplicates overlapping items across sources,
\item
  stores recurring findings in SQLite,
\item
  supports a watchlist and repeated scheduled reruns,
\item
  renders compact, grounded briefings with citations.
\end{itemize}

Its core architecture is:

\begin{Shaded}
\begin{Highlighting}[]
\NormalTok{last30days.py orchestrator}
\NormalTok{  {-}\textgreater{} source connectors}
\NormalTok{  {-}\textgreater{} normalization}
\NormalTok{  {-}\textgreater{} scoring}
\NormalTok{  {-}\textgreater{} dedupe}
\NormalTok{  {-}\textgreater{} rendering}
\NormalTok{  {-}\textgreater{} optional watchlist/store}
\end{Highlighting}
\end{Shaded}

The repo is methodologically stronger than it first appears because it
encodes:

\begin{itemize}
\tightlist
\item
  query expansion,
\item
  source fusion,
\item
  canonical schemas,
\item
  confidence-aware scoring,
\item
  persistent research memory,
\item
  repeated watchlist execution,
\item
  report generation meant for human decision support.
\end{itemize}

\subsubsection{0.3 How the reference repo applies to your
project}\label{03-how-the-reference-repo-applies-to-your-project}

It is valuable less as direct market-data code and more as an
architecture pattern for the intelligence layer you want to build on top
of your collector.

Direct benefits to your project:

\begin{itemize}
\tightlist
\item
  It demonstrates how to separate ingestion, normalization, scoring,
  storage, and rendering.
\item
  It gives you a working pattern for a watchlist and recurring topic
  research loop.
\item
  It shows how to fuse structured signals and unstructured external
  evidence into a single briefing.
\item
  It treats Polymarket as one source among many, which is exactly how
  your alerting layer should behave once a market anomaly is detected.
\item
  Its \texttt{store.py} and \texttt{watchlist.py} are a good template
  for persistent signal history, briefing history, and scheduled reruns.
\item
  Its scoring stack is a useful conceptual starting point for ranking
  candidate signals by evidence quality instead of raw magnitude alone.
\end{itemize}

\subsubsection{0.4 What you can reuse directly vs. what needs
adaptation}\label{04-what-you-can-reuse-directly-vs-what-needs-adaptation}

Directly reusable patterns:

\begin{itemize}
\tightlist
\item
  Watchlist and recurring-run database pattern from
  \texttt{scripts/store.py}.
\item
  CLI task routing pattern from \texttt{scripts/watchlist.py}.
\item
  Canonical-schema normalization mindset from
  \texttt{scripts/lib/normalize.py}.
\item
  Scoring pipeline structure from \texttt{scripts/lib/score.py}.
\item
  Dedupe and cross-reference pattern.
\item
  Rendering and compact briefing pattern.
\item
  Query expansion and topic enrichment pattern.
\end{itemize}

Reusable with adaptation:

\begin{itemize}
\tightlist
\item
  Polymarket connector ideas from \texttt{scripts/lib/polymarket.py}.
\item
  Evidence fusion across market + X + web.
\item
  Research-memory SQLite schema.
\item
  Alert ranking logic.
\end{itemize}

Needs major adaptation:

\begin{itemize}
\tightlist
\item
  The reference repo is topic-centric; your system must be
  market-centric, wallet-centric, and event-timestamp-centric.
\item
  Its scoring is built for content relevance and engagement; yours must
  score microstructure anomalies, order-flow coordination, and
  predictive lead.
\item
  Its data model is item-based; yours needs entities for wallets,
  clusters, episodes, event windows, and alerts.
\item
  Its storage is good for research notes; yours needs durable
  time-series and account-level state.
\item
  Its Polymarket usage is search-level and descriptive; yours needs
  continuous ingestion and position-level inference.
\end{itemize}

Not worth reusing directly:

\begin{itemize}
\tightlist
\item
  The skill packaging, first-run setup flow, and user-facing assistant
  UX scaffolding.
\item
  Most platform-specific social-search clients unless you want the exact
  same providers.
\item
  Generic prompt-research heuristics that are not specific to market
  intelligence.
\end{itemize}

\subsubsection{0.5 Practical synthesis}\label{05-practical-synthesis}

Your repo should remain the canonical market-data plane.

The professor repo should inspire the evidence-fusion and briefing
plane.

The most productive mental model is:

\begin{Shaded}
\begin{Highlighting}[]
\NormalTok{Current repo = sensing layer for Polymarket microstructure}
\NormalTok{Reference repo = intelligence/reporting layer for external context fusion}
\NormalTok{Target system = real{-}time anomaly + evidence fusion + alerting platform}
\end{Highlighting}
\end{Shaded}

\subsection{1. Executive Summary}\label{1-executive-summary}

This project should become a real-time information-crunching engine for
prediction markets, designed to detect when market prices are moving for
reasons that appear better informed than the public narrative.

The core research contribution is not simply "predict Polymarket." It is
to estimate whether specific kinds of order flow, especially
concentrated activity from apparently new or previously inactive
wallets, can identify information arrival earlier than conventional
public signals.

The strongest research framing is:

\begin{itemize}
\tightlist
\item
  detect anomalous coordinated order flow,
\item
  estimate whether it is informed rather than noisy,
\item
  test whether it leads public information discovery,
\item
  and convert the result into ranked, explainable alerts.
\end{itemize}

This is potentially novel because it combines:

\begin{itemize}
\tightlist
\item
  prediction market microstructure,
\item
  wallet-level behavioral signatures,
\item
  external evidence retrieval,
\item
  and tradable event-study validation.
\end{itemize}

It also has a built-in research tension that makes it stronger:

\begin{itemize}
\tightlist
\item
  if the signal becomes widely watched, its alpha may decay,
\item
  but its calibration or social usefulness may improve.
\end{itemize}

That tension is the right way to frame your Goodhart question.

\subsection{2. System Overview}\label{2-system-overview}

\subsubsection{2.1 Target architecture}\label{21-target-architecture}

\begin{Shaded}
\begin{Highlighting}[]
\NormalTok{                +{-}{-}{-}{-}{-}{-}{-}{-}{-}{-}{-}{-}{-}{-}{-}{-}{-}{-}{-}{-}{-}{-}+}
\NormalTok{                |  Polymarket APIs     |}
\NormalTok{                |  Gamma / Data / WSS  |}
\NormalTok{                +{-}{-}{-}{-}{-}{-}{-}{-}{-}{-}+{-}{-}{-}{-}{-}{-}{-}{-}{-}{-}{-}+}
\NormalTok{                           |}
\NormalTok{                           v}
\NormalTok{              +{-}{-}{-}{-}{-}{-}{-}{-}{-}{-}{-}{-}{-}{-}{-}{-}{-}{-}{-}{-}{-}{-}{-}{-}{-}{-}{-}+}
\NormalTok{              |  Ingestion Layer          |}
\NormalTok{              |  {-} market sync            |}
\NormalTok{              |  {-} websocket listener     |}
\NormalTok{              |  {-} trades/activity pull   |}
\NormalTok{              |  {-} positions pull         |}
\NormalTok{              +{-}{-}{-}{-}{-}{-}{-}{-}{-}{-}{-}{-}+{-}{-}{-}{-}{-}{-}{-}{-}{-}{-}{-}{-}{-}{-}+}
\NormalTok{                           |}
\NormalTok{                           v}
\NormalTok{              +{-}{-}{-}{-}{-}{-}{-}{-}{-}{-}{-}{-}{-}{-}{-}{-}{-}{-}{-}{-}{-}{-}{-}{-}{-}{-}{-}+}
\NormalTok{              |  Storage Layer            |}
\NormalTok{              |  {-} Postgres/Timeseries    |}
\NormalTok{              |  {-} raw events             |}
\NormalTok{              |  {-} normalized features    |}
\NormalTok{              |  {-} wallet state           |}
\NormalTok{              |  {-} signal history         |}
\NormalTok{              +{-}{-}{-}{-}{-}{-}{-}{-}{-}{-}{-}{-}+{-}{-}{-}{-}{-}{-}{-}{-}{-}{-}{-}{-}{-}{-}+}
\NormalTok{                           |}
\NormalTok{                           v}
\NormalTok{              +{-}{-}{-}{-}{-}{-}{-}{-}{-}{-}{-}{-}{-}{-}{-}{-}{-}{-}{-}{-}{-}{-}{-}{-}{-}{-}{-}+}
\NormalTok{              |  Detection Layer          |}
\NormalTok{              |  {-} new wallet detector    |}
\NormalTok{              |  {-} coordinated{-}flow model |}
\NormalTok{              |  {-} price/CI state model   |}
\NormalTok{              |  {-} anomaly scoring        |}
\NormalTok{              +{-}{-}{-}{-}{-}{-}{-}{-}{-}{-}{-}{-}+{-}{-}{-}{-}{-}{-}{-}{-}{-}{-}{-}{-}{-}{-}+}
\NormalTok{                           |}
\NormalTok{                           v}
\NormalTok{              +{-}{-}{-}{-}{-}{-}{-}{-}{-}{-}{-}{-}{-}{-}{-}{-}{-}{-}{-}{-}{-}{-}{-}{-}{-}{-}{-}+}
\NormalTok{              |  Evidence Layer           |}
\NormalTok{              |  {-} X recent search        |}
\NormalTok{              |  {-} web/news search        |}
\NormalTok{              |  {-} market correlation     |}
\NormalTok{              |  {-} narrative extraction   |}
\NormalTok{              +{-}{-}{-}{-}{-}{-}{-}{-}{-}{-}{-}{-}+{-}{-}{-}{-}{-}{-}{-}{-}{-}{-}{-}{-}{-}{-}+}
\NormalTok{                           |}
\NormalTok{                           v}
\NormalTok{              +{-}{-}{-}{-}{-}{-}{-}{-}{-}{-}{-}{-}{-}{-}{-}{-}{-}{-}{-}{-}{-}{-}{-}{-}{-}{-}{-}+}
\NormalTok{              |  Alerting Layer           |}
\NormalTok{              |  {-} rank and suppress      |}
\NormalTok{              |  {-} Telegram/Discord/SMS   |}
\NormalTok{              |  {-} analyst dashboard      |}
\NormalTok{              +{-}{-}{-}{-}{-}{-}{-}{-}{-}{-}{-}{-}+{-}{-}{-}{-}{-}{-}{-}{-}{-}{-}{-}{-}{-}{-}+}
\NormalTok{                           |}
\NormalTok{                           v}
\NormalTok{              +{-}{-}{-}{-}{-}{-}{-}{-}{-}{-}{-}{-}{-}{-}{-}{-}{-}{-}{-}{-}{-}{-}{-}{-}{-}{-}{-}+}
\NormalTok{              |  Research Layer           |}
\NormalTok{              |  {-} backtests              |}
\NormalTok{              |  {-} event studies          |}
\NormalTok{              |  {-} model training         |}
\NormalTok{              |  {-} paper{-}trading          |}
\NormalTok{              +{-}{-}{-}{-}{-}{-}{-}{-}{-}{-}{-}{-}{-}{-}{-}{-}{-}{-}{-}{-}{-}{-}{-}{-}{-}{-}{-}+}
\end{Highlighting}
\end{Shaded}

\subsubsection{2.2 Component map}\label{22-component-map}

\begin{enumerate}
\def\labelenumi{\arabic{enumi}.}
\item
  Market ingestion service

  \begin{itemize}
  \tightlist
  \item
    Pulls Gamma events and markets.
  \item
    Maintains websocket market subscriptions.
  \item
    Pulls public trades, user activity, and market position snapshots
    for flagged markets.
  \end{itemize}
\item
  Normalization and feature service

  \begin{itemize}
  \tightlist
  \item
    Converts raw payloads into market, wallet, cluster, and episode
    tables.
  \item
    Computes rolling features.
  \end{itemize}
\item
  Detection engine

  \begin{itemize}
  \tightlist
  \item
    Produces anomaly scores, leadership scores, and signal confidence.
  \end{itemize}
\item
  External evidence engine

  \begin{itemize}
  \tightlist
  \item
    Searches X and web/news after a candidate anomaly.
  \item
    Summarizes whether public explanation exists already.
  \end{itemize}
\item
  Alerting and analyst interface

  \begin{itemize}
  \tightlist
  \item
    Delivers ranked alerts.
  \item
    Stores alert outcomes and analyst feedback.
  \end{itemize}
\item
  Backtesting and experimentation environment

  \begin{itemize}
  \tightlist
  \item
    Reconstructs point-in-time states.
  \item
    Evaluates predictive, statistical, and economic performance.
  \end{itemize}
\end{enumerate}

\subsubsection{2.3 Core entity model}\label{23-core-entity-model}

Required primary entities:

\begin{itemize}
\tightlist
\item
  \texttt{events}
\item
  \texttt{markets}
\item
  \texttt{market\_snapshots}
\item
  \texttt{order\_book\_snapshots}
\item
  \texttt{trades\_raw}
\item
  \texttt{wallets}
\item
  \texttt{wallet\_profiles}
\item
  \texttt{wallet\_market\_positions}
\item
  \texttt{wallet\_activity}
\item
  \texttt{wallet\_clusters}
\item
  \texttt{signal\_episodes}
\item
  \texttt{external\_evidence}
\item
  \texttt{alerts}
\item
  \texttt{alert\_outcomes}
\end{itemize}

\subsection{3. Research Hypotheses}\label{3-research-hypotheses}

\subsubsection{3.1 Terminology
correction}\label{31-terminology-correction}

Your phrase "confidence intervals rise" needs tightening.

If you mean a statistical confidence interval, a wider interval means
more uncertainty, not stronger conviction. The system should separately
track:

\begin{itemize}
\tightlist
\item
  \texttt{p\_t}: market-implied probability,
\item
  \texttt{u\_t}: uncertainty or credible-interval width,
\item
  \texttt{c\_t}: confidence score derived from liquidity, order-flow
  consistency, and external corroboration.
\end{itemize}

The economically interesting pattern is usually:

\begin{itemize}
\tightlist
\item
  a fast move in \texttt{p\_t},
\item
  accompanied by strong or coordinated order flow,
\item
  with stable or narrowing \texttt{u\_t},
\item
  and rising \texttt{c\_t}.
\end{itemize}

\subsubsection{3.2 H1: Goodhart\textquotesingle s Law immunity
test}\label{32-h1-goodharts-law-immunity-test}

Research question:

Does a Polymarket-derived signal lose predictive or tradable value once
it becomes widely observed and operationalized by others?

Null hypothesis H1\_0:

Conditional on event type and liquidity, wider public observability of a
market signal does not change its predictive value, calibration, or
post-signal alpha.

Alternative hypothesis H1\_A:

As observability rises, at least one of the following deteriorates:

\begin{itemize}
\tightlist
\item
  predictive lead over public information,
\item
  post-signal alpha,
\item
  calibration quality,
\item
  or signal-to-noise ratio.
\end{itemize}

Critical challenge:

Goodhart may not show up as worse calibration. It may show up as lower
tradable edge because the market becomes more efficient after being
watched. So this test must separate:

\begin{itemize}
\tightlist
\item
  forecast quality,
\item
  market quality,
\item
  and extractable alpha.
\end{itemize}

Recommended test design:

\begin{enumerate}
\def\labelenumi{\arabic{enumi}.}
\item
  Define observability shocks.

  \begin{itemize}
  \tightlist
  \item
    Candidate shocks:

    \begin{itemize}
    \tightlist
    \item
      market embedded in widely viewed dashboards,
    \item
      large media citation volume,
    \item
      repeated mentions by major X accounts,
    \item
      internal alert launch date,
    \item
      public leaderboard/dashboard publication.
    \end{itemize}
  \end{itemize}
\item
  Build matched treatment and control groups.

  \begin{itemize}
  \tightlist
  \item
    Match on category, liquidity, age, resolution horizon, and baseline
    volatility.
  \end{itemize}
\item
  Measure pre/post differences in:

  \begin{itemize}
  \tightlist
  \item
    Brier score and log loss at fixed horizons,
  \item
    price impact of new public information,
  \item
    spread and depth,
  \item
    order-flow toxicity,
  \item
    abnormal returns after your signal,
  \item
    and alert precision.
  \end{itemize}
\item
  Use a difference-in-differences or staggered event-study design.
\item
  Run placebo shocks on unaffected markets.
\end{enumerate}

Success criterion:

If calibration improves but alpha decays, the market is not "immune"; it
is becoming more efficient under observation.

\subsubsection{3.3 H2: Insider activity detection via new wallet
signatures}\label{33-h2-insider-activity-detection-via-new-wallet-signatures}

Research question:

Do newly observed wallets, especially those making large first trades
and converging on the same event, contain predictive information about
near-term event repricing?

Null hypothesis H2\_0:

New-wallet large-trade flow has no incremental predictive value after
controlling for price momentum, liquidity, and generic order imbalance.

Alternative hypothesis H2\_A:

Coordinated new-wallet flow predicts subsequent repricing or
public-information arrival beyond baseline market microstructure
features.

Operational definitions:

\begin{itemize}
\tightlist
\item
  New wallet:

  \begin{itemize}
  \tightlist
  \item
    first appearance in your local history within lookback window
    \texttt{L},
  \item
    and low or zero prior Polymarket footprint based on:

    \begin{itemize}
    \tightlist
    \item
      no observed local trades,
    \item
      optional \texttt{traded} count near zero,
    \item
      minimal public profile history if available.
    \end{itemize}
  \end{itemize}
\item
  Large first trade:

  \begin{itemize}
  \tightlist
  \item
    first trade notional greater than \texttt{\$10,000},
  \item
    plus category-adjusted threshold such as
    \texttt{\textgreater{}\ p95} of first-trade size in that market
    bucket.
  \end{itemize}
\item
  Coordinated flow episode:

  \begin{itemize}
  \tightlist
  \item
    at least \texttt{N} new wallets,
  \item
    on same condition or event family,
  \item
    within time window \texttt{T},
  \item
    showing same directional exposure.
  \end{itemize}
\end{itemize}

Critical challenge:

New proxy wallet is not the same as new human. Polymarket supports
proxy/safe wallet structures, so wallet novelty is only a behavioral
proxy. This project must never equate "new wallet" with "new insider."
It should use terms such as:

\begin{itemize}
\tightlist
\item
  anomalous wallet,
\item
  previously unseen wallet,
\item
  coordinated fresh-wallet flow,
\item
  potential informed flow.
\end{itemize}

Required detection features:

\begin{itemize}
\tightlist
\item
  first-trade notional,
\item
  cumulative notional over 5m, 15m, 1h,
\item
  number of fresh wallets entering event,
\item
  concentration ratio of top fresh wallets,
\item
  directional order imbalance,
\item
  share of market volume from fresh wallets,
\item
  sequence density and synchronization,
\item
  persistence of position growth after first trade,
\item
  overlap across related markets,
\item
  external-public-silence score.
\end{itemize}

Key public endpoints that make H2 feasible:

\begin{itemize}
\tightlist
\item
  \texttt{/trades} returns \texttt{proxyWallet}, \texttt{conditionId},
  \texttt{asset}, \texttt{price}, \texttt{size}, \texttt{timestamp},
  \texttt{transactionHash}.
\item
  \texttt{/activity} returns \texttt{proxyWallet}, \texttt{type},
  \texttt{size}, \texttt{usdcSize}, \texttt{price}, \texttt{side}, and
  \texttt{conditionId}.
\item
  \texttt{/positions} and \texttt{/positions?market=...} support
  checking current wallet exposure.
\item
  \texttt{/public-profile} exposes public profile metadata and profile
  creation information when available.
\item
  \texttt{/traded} gives total markets traded by a user.
\end{itemize}

\subsubsection{3.4 H3: Leading indicator
construction}\label{34-h3-leading-indicator-construction}

Research question:

Can insider-activity-weighted Polymarket signals lead real-world events
or lead public recognition of those events?

Null hypothesis H3\_0:

Signals built from fresh-wallet and coordinated-flow features do not
forecast subsequent public information arrival or profitable market
repricing better than naive Polymarket momentum.

Alternative hypothesis H3\_A:

Signals built from wallet-aware order flow lead:

\begin{itemize}
\tightlist
\item
  future Polymarket repricing,
\item
  cross-market spillovers,
\item
  or external-public evidence arrival,
\end{itemize}

with meaningful economic value.

Validation targets:

\begin{itemize}
\tightlist
\item
  time to first major public corroboration,
\item
  time to cross-market repricing,
\item
  time to realized event occurrence,
\item
  short-horizon abnormal returns after signal,
\item
  alert precision at top-k.
\end{itemize}

\subsubsection{3.5 Additional hypotheses worth
testing}\label{35-additional-hypotheses-worth-testing}

H4: Order-flow toxicity hypothesis

Markets with abnormal buy-side imbalance from large trades and shrinking
passive depth exhibit larger next-interval price moves than markets with
similar momentum but healthier depth.

H5: Cross-market leakage hypothesis

Information often enters narrower or derivative markets before repricing
broader headline markets. For example, related conditional or sector
markets may move before the flagship event contract.

H6: Public-silence hypothesis

An anomaly is more predictive when market flow intensifies before
comparable growth in X/news discussion.

H7: Attention-alpha tradeoff hypothesis

Higher market observability improves liquidity and calibration but
reduces post-alert alpha.

H8: Yes-bias correction hypothesis

Prediction markets may exhibit systematic default or "Yes" overtrading
in certain contexts; signals that adjust for this behavioral bias will
outperform raw flow signals.

\subsubsection{3.6 Where classic quantitative finance applies and
fails}\label{36-where-classic-quantitative-finance-applies-and-fails}

Applies well:

\begin{itemize}
\tightlist
\item
  market microstructure,
\item
  order imbalance,
\item
  price impact,
\item
  point-process models,
\item
  state-space filtering,
\item
  change-point detection,
\item
  toxicity metrics such as PIN/VPIN-style ideas,
\item
  event studies,
\item
  lead-lag analysis,
\item
  graph clustering of flow.
\end{itemize}

Applies with adaptation:

\begin{itemize}
\tightlist
\item
  factor-style decomposition across event families,
\item
  latent-state models for macro regime and event urgency,
\item
  cross-sectional ranking models.
\end{itemize}

Fails or weakens:

\begin{itemize}
\tightlist
\item
  classic discounted-cash-flow intuition,
\item
  standard equity factor models built on firm fundamentals,
\item
  assumptions of stationary return distributions,
\item
  assumptions of continuous unbounded prices,
\item
  assumptions that contracts represent the same economic claim over long
  windows.
\end{itemize}

Why:

\begin{itemize}
\tightlist
\item
  prices are bounded between 0 and 1,
\item
  resolution causes hard terminal jumps,
\item
  contract lifetimes are finite and heterogeneous,
\item
  market creation/resolution rules matter,
\item
  and exogenous event arrival dominates many dynamics.
\end{itemize}

\subsection{4. Data Pipeline
Requirements}\label{4-data-pipeline-requirements}

\subsubsection{4.1 Current local pipeline
limitations}\label{41-current-local-pipeline-limitations}

Current limitations that must be fixed before serious modeling:

\begin{enumerate}
\def\labelenumi{\arabic{enumi}.}
\item
  Coverage mismatch

  \begin{itemize}
  \tightlist
  \item
    Sports-only filtering currently excludes much of the research target
    universe.
  \end{itemize}
\item
  Identity blindness

  \begin{itemize}
  \tightlist
  \item
    No wallet/account tables.
  \item
    No profile metadata.
  \item
    No market-position snapshots by wallet.
  \end{itemize}
\item
  Incomplete trade capture

  \begin{itemize}
  \tightlist
  \item
    REST trade polling/backfill appears misconfigured against the
    documented API contract.
  \end{itemize}
\item
  One-sided asset coverage

  \begin{itemize}
  \tightlist
  \item
    Only yes-token subscriptions are used in the live path.
  \end{itemize}
\item
  Research-hostile retention

  \begin{itemize}
  \tightlist
  \item
    Closed-market detailed data is removed after 24 hours.
  \end{itemize}
\item
  Storage bottleneck

  \begin{itemize}
  \tightlist
  \item
    Single 11 GB SQLite database is fine for local experimentation but
    not for 24/7 multi-service analytics.
  \end{itemize}
\item
  No queueing or replay

  \begin{itemize}
  \tightlist
  \item
    Raw events are not durably queued for reprocessing.
  \end{itemize}
\item
  No quality-control tables

  \begin{itemize}
  \tightlist
  \item
    No ingestion lag metrics, dead-letter queues, schema versioning
    tables, or alert audit tables.
  \end{itemize}
\end{enumerate}

\subsubsection{4.2 Required pipeline
upgrades}\label{42-required-pipeline-upgrades}

Must-have ingestion changes:

\begin{itemize}
\tightlist
\item
  Remove hard sports filter or make it configurable.
\item
  Subscribe to both yes and no assets where supported by market
  structure.
\item
  Store raw trades with:

  \begin{itemize}
  \tightlist
  \item
    \texttt{proxy\_wallet}
  \item
    \texttt{transaction\_hash}
  \item
    \texttt{condition\_id}
  \item
    \texttt{asset}
  \item
    \texttt{outcome\_side}
  \item
    \texttt{price}
  \item
    \texttt{size}
  \item
    \texttt{usdc\_notional}
  \item
    \texttt{timestamp}
  \item
    \texttt{source}
  \end{itemize}
\item
  Add wallet profile enrichment from public profile endpoints.
\item
  Add wallet activity and position snapshot pulls for flagged markets.
\item
  Preserve closed-market microstructure for at least 30-90 days in cloud
  cold storage.
\item
  Write raw payloads to append-only storage before normalization.
\end{itemize}

Required schema additions:

\begin{itemize}
\tightlist
\item
  \texttt{wallets}
\item
  \texttt{wallet\_profiles}
\item
  \texttt{wallet\_first\_seen}
\item
  \texttt{wallet\_activity}
\item
  \texttt{wallet\_positions}
\item
  \texttt{wallet\_market\_exposure\_rollups}
\item
  \texttt{signal\_candidates}
\item
  \texttt{signal\_features}
\item
  \texttt{signals\_validated}
\item
  \texttt{alerts\_sent}
\end{itemize}

\subsubsection{4.3 Suggested storage
model}\label{43-suggested-storage-model}

Recommended split:

\begin{itemize}
\tightlist
\item
  Postgres for relational state and alerting.
\item
  Time-partitioned tables for trades/snapshots/activity.
\item
  Object storage for raw payload archives and backtest snapshots.
\item
  Redis for short-term caches, queues, and suppression windows.
\end{itemize}

Recommended retention:

\begin{itemize}
\tightlist
\item
  raw websocket/trade payloads: 30-90 days hot, 1+ year cold
\item
  normalized trades/activity: 1+ year
\item
  feature tables: 1+ year
\item
  alert outcomes and labels: permanent
\end{itemize}

\subsubsection{4.4 Cloud migration
recommendation}\label{44-cloud-migration-recommendation}

Best choice among your listed options:

\textbf{GCP Cloud Run} for the first serious 24/7 version.

Why Cloud Run wins here:

\begin{itemize}
\tightlist
\item
  Easier than AWS to stand up quickly.
\item
  Supports WebSockets.
\item
  Supports minimum instances for warm always-on services.
\item
  Works well with Cloud Scheduler for recurring jobs.
\item
  Clean integration with managed Postgres and object storage.
\end{itemize}

Important caveat:

Cloud Run WebSockets are still HTTP requests and remain subject to
request timeout. Google documents a default timeout of 5 minutes and a
maximum under 60 minutes, so long-lived stream workers must reconnect by
design.

Recommended GCP deployment:

\begin{itemize}
\tightlist
\item
  Cloud Run service \texttt{collector-realtime}

  \begin{itemize}
  \tightlist
  \item
    websocket listener plus lightweight pollers
  \end{itemize}
\item
  Cloud Run service \texttt{collector-batch}

  \begin{itemize}
  \tightlist
  \item
    metadata sync, wallet enrichment, backfills
  \end{itemize}
\item
  Cloud Run service \texttt{detector}

  \begin{itemize}
  \tightlist
  \item
    signal scoring and alert generation
  \end{itemize}
\item
  Cloud Scheduler

  \begin{itemize}
  \tightlist
  \item
    periodic sync, integrity checks, retraining triggers
  \end{itemize}
\item
  Cloud SQL Postgres

  \begin{itemize}
  \tightlist
  \item
    normalized operational store
  \end{itemize}
\item
  Cloud Storage

  \begin{itemize}
  \tightlist
  \item
    raw payload archive
  \end{itemize}
\item
  Redis or Memorystore

  \begin{itemize}
  \tightlist
  \item
    caches, locks, suppression windows
  \end{itemize}
\end{itemize}

Option tradeoffs:

\textbf{AWS Lambda + SQS}

Pros:

\begin{itemize}
\tightlist
\item
  Excellent for decoupled async workers.
\item
  Strong retry and scaling story.
\item
  Good for enrichment, feature jobs, and alert fan-out.
\end{itemize}

Cons:

\begin{itemize}
\tightlist
\item
  Poor fit for a persistent websocket collector.
\item
  Lambda is stateless and time-bounded, so continuous market listening
  becomes awkward.
\item
  If you stay on AWS, ECS/Fargate is a better core-ingestion choice than
  Lambda.
\end{itemize}

Recommendation:

\begin{itemize}
\tightlist
\item
  Use Lambda + SQS only for auxiliary workers, not the primary realtime
  collector.
\end{itemize}

\textbf{GCP Cloud Run}

Pros:

\begin{itemize}
\tightlist
\item
  Best balance of ease, cost, and operational maturity for a research
  platform.
\item
  Native cron-like scheduling via Cloud Scheduler.
\item
  Supports min instances and WebSockets.
\end{itemize}

Cons:

\begin{itemize}
\tightlist
\item
  WebSocket sessions must reconnect because of request timeout limits.
\item
  Cost management requires attention if minimum instances are always
  warm.
\end{itemize}

Recommendation:

\begin{itemize}
\tightlist
\item
  Best overall choice among the options you listed.
\end{itemize}

\textbf{Railway}

Pros:

\begin{itemize}
\tightlist
\item
  Fastest path to a working MVP.
\item
  Friendly for a single persistent service and quick iteration.
\item
  Supports volumes for persistent local-like storage.
\end{itemize}

Cons:

\begin{itemize}
\tightlist
\item
  Cron jobs are explicitly for short-lived tasks, not long-running
  workers.
\item
  Volume limitations and redeploy behavior are more constraining for a
  serious data platform.
\item
  Better for prototype and demo than a thesis-grade always-on research
  backbone.
\end{itemize}

Recommendation:

\begin{itemize}
\tightlist
\item
  Use for MVP only if speed matters more than rigor.
\end{itemize}

\subsection{5. Real-Time Detection
Engine}\label{5-real-time-detection-engine}

\subsubsection{5.1 Design principle}\label{51-design-principle}

The detection engine should be multi-stage, not monolithic.

Do not let one model directly emit final alerts from raw ticks.

Instead use a staged pipeline:

\begin{enumerate}
\def\labelenumi{\arabic{enumi}.}
\tightlist
\item
  candidate generation
\item
  wallet-flow aggregation
\item
  anomaly scoring
\item
  external-evidence check
\item
  alert ranking and suppression
\end{enumerate}

\subsubsection{5.2 Agentic workflow
specification}\label{52-agentic-workflow-specification}

Use specialized services or "agents," but keep internal reasoning
private. The system should store evidence, scores, and concise
rationales, not raw hidden chain-of-thought.

Suggested agents:

\begin{enumerate}
\def\labelenumi{\arabic{enumi}.}
\item
  Market anomaly agent

  \begin{itemize}
  \tightlist
  \item
    Watches repricing speed, spread changes, depth imbalance, and
    volatility jumps.
  \end{itemize}
\item
  Wallet-flow agent

  \begin{itemize}
  \tightlist
  \item
    Tracks fresh wallets, first-trade notional, concentration,
    clustering, and persistence.
  \end{itemize}
\item
  Market-linkage agent

  \begin{itemize}
  \tightlist
  \item
    Checks neighboring markets, same event tree, and related condition
    IDs.
  \end{itemize}
\item
  External-evidence agent

  \begin{itemize}
  \tightlist
  \item
    Queries X and web/news after a candidate anomaly.
  \item
    Distinguishes:

    \begin{itemize}
    \tightlist
    \item
      public explanation already visible,
    \item
      weak or ambiguous public chatter,
    \item
      no visible public explanation.
    \end{itemize}
  \end{itemize}
\item
  Alert-composer agent

  \begin{itemize}
  \tightlist
  \item
    Produces human-readable summary with ranked evidence and
    invalidation conditions.
  \end{itemize}
\end{enumerate}

\subsubsection{5.3 X and web/Google
integration}\label{53-x-and-webgoogle-integration}

Recommended implementation:

\begin{itemize}
\tightlist
\item
  X Recent Search API for rapid context windows.
\item
  Web/news search provider for broader corroboration.
\end{itemize}

Important current constraint:

\begin{itemize}
\tightlist
\item
  X\textquotesingle s documented recent-search endpoint retrieves posts
  from the last 7 days and supports up to 100 results per request. It is
  good for near-term corroboration, not deep history.
\item
  Google\textquotesingle s Custom Search JSON API is closed to new
  customers and existing customers must transition by January 1, 2027.
  Do not design a new core dependency around it.
\end{itemize}

Practical recommendation:

\begin{itemize}
\tightlist
\item
  Use X Recent Search for fast corroboration.
\item
  Use a general web/news search provider for full-web evidence
  retrieval.
\item
  If you specifically need Google-style search results, treat that as a
  provider abstraction, not a hard-coded dependency.
\end{itemize}

\subsubsection{5.4 Candidate signal
logic}\label{54-candidate-signal-logic}

A candidate alert episode should require:

\begin{itemize}
\tightlist
\item
  unusual probability velocity,
\item
  abnormal fresh-wallet flow,
\item
  or abnormal large-trade imbalance,
\item
  plus at least one supporting context feature.
\end{itemize}

Example candidate rule:

\begin{Shaded}
\begin{Highlighting}[]
\NormalTok{Episode triggers when:}
\NormalTok{  fresh\_wallet\_notional\_15m \textgreater{} threshold\_A}
\NormalTok{  and fresh\_wallet\_count\_15m \textgreater{} threshold\_B}
\NormalTok{  and probability\_change\_15m \textgreater{} threshold\_C}
\NormalTok{  and external\_public\_explanation\_score \textless{} threshold\_D}
\end{Highlighting}
\end{Shaded}

\subsubsection{5.5 Core signal features}\label{55-core-signal-features}

Market-state features:

\begin{itemize}
\tightlist
\item
  implied probability change
\item
  bid-ask spread
\item
  depth imbalance
\item
  order-book resiliency
\item
  realized volatility
\item
  volume acceleration
\end{itemize}

Wallet features:

\begin{itemize}
\tightlist
\item
  first seen timestamp
\item
  first trade size
\item
  number of markets traded historically
\item
  public profile age if available
\item
  concentration by wallet and cluster
\item
  repeated directional reinforcement
\end{itemize}

Cross-context features:

\begin{itemize}
\tightlist
\item
  correlated market confirmation
\item
  event family breadth
\item
  X/news corroboration lag
\item
  contradiction score from public sources
\end{itemize}

\subsubsection{5.6 Signal output}\label{56-signal-output}

Each signal should emit:

\begin{itemize}
\tightlist
\item
  \texttt{signal\_id}
\item
  event and market identifiers
\item
  direction
\item
  anomaly score
\item
  confidence score
\item
  supporting wallets and concentration stats
\item
  market microstructure summary
\item
  external evidence summary
\item
  invalidation conditions
\end{itemize}

\subsection{6. Alert and Notification
System}\label{6-alert-and-notification-system}

\subsubsection{6.1 Channel evaluation}\label{61-channel-evaluation}

\textbf{Telegram Bot API}

Pros:

\begin{itemize}
\tightlist
\item
  Free.
\item
  HTTP-based and easy to automate.
\item
  Strong mobile push experience.
\item
  Supports rich text, inline buttons, and media.
\end{itemize}

Cons:

\begin{itemize}
\tightlist
\item
  Consumer rather than enterprise audit/compliance tooling.
\item
  Requires user/bot setup.
\end{itemize}

Best use:

\begin{itemize}
\tightlist
\item
  Primary 24/7 personal alert hotline.
\end{itemize}

\textbf{Discord Webhooks}

Pros:

\begin{itemize}
\tightlist
\item
  Very low effort.
\item
  Easy structured embeds.
\item
  Good for team rooms and historical channels.
\end{itemize}

Cons:

\begin{itemize}
\tightlist
\item
  Better for stream-of-alerts than high-urgency personal wake-up alerts.
\item
  Easy for important alerts to get buried in channel noise.
\end{itemize}

Best use:

\begin{itemize}
\tightlist
\item
  Secondary team channel, archive sink, and debugging channel.
\end{itemize}

\textbf{Twilio SMS}

Pros:

\begin{itemize}
\tightlist
\item
  Highest urgency.
\item
  Works even when app push is muted or unavailable.
\item
  Message lifecycle and delivery states are mature.
\end{itemize}

Cons:

\begin{itemize}
\tightlist
\item
  Direct cost.
\item
  Strictest anti-spam and phone-number operational burden.
\item
  Too noisy if used for anything but highest-severity alerts.
\end{itemize}

Best use:

\begin{itemize}
\tightlist
\item
  Escalation path only for the top few alerts per week or for repeated
  confirmations.
\end{itemize}

\textbf{ntfy.sh}

Pros:

\begin{itemize}
\tightlist
\item
  Extremely simple HTTP publish model.
\item
  Fast for personal prototypes.
\item
  Good fallback or self-hostable option.
\end{itemize}

Cons:

\begin{itemize}
\tightlist
\item
  Less institutional and less universal than Telegram/SMS.
\item
  Topic secrecy is part of the security model.
\end{itemize}

Best use:

\begin{itemize}
\tightlist
\item
  Developer notifications, laptop/mobile fallback, or lab use.
\end{itemize}

\subsubsection{6.2 Recommended alert
stack}\label{62-recommended-alert-stack}

Recommended:

\begin{itemize}
\tightlist
\item
  Primary: Telegram
\item
  Secondary: Discord
\item
  Escalation: Twilio SMS
\item
  Developer fallback: ntfy
\end{itemize}

Why:

Telegram gives the best balance of cost, reliability, message richness,
and mobile usability for a single researcher running 24/7 monitoring.

\subsubsection{6.3 Alert severity
levels}\label{63-alert-severity-levels}

\begin{itemize}
\tightlist
\item
  \texttt{INFO}: interesting but non-actionable
\item
  \texttt{WATCH}: monitor closely
\item
  \texttt{ACTIONABLE}: high-confidence signal worth checking immediately
\item
  \texttt{CRITICAL}: repeated or multi-source-confirmed high-value
  signal
\end{itemize}

\subsubsection{6.4 Alert schema: top 5 things you need to
know}\label{64-alert-schema-top-5-things-you-need-to-know}

Every alert should surface these five items first:

\begin{enumerate}
\def\labelenumi{\arabic{enumi}.}
\item
  What changed

  \begin{itemize}
  \tightlist
  \item
    event, market, direction, probability move, and time window
  \end{itemize}
\item
  Why it looks informed

  \begin{itemize}
  \tightlist
  \item
    fresh-wallet count, first-trade sizes, concentration, coordinated
    directionality
  \end{itemize}
\item
  How strong the market evidence is

  \begin{itemize}
  \tightlist
  \item
    anomaly score, confidence score, volume share, spread/depth
    conditions
  \end{itemize}
\item
  What public evidence exists already

  \begin{itemize}
  \tightlist
  \item
    top X/news hits, first-public-mention lag, contradictions,
    unresolved catalyst
  \end{itemize}
\item
  What action framework follows

  \begin{itemize}
  \tightlist
  \item
    entry thesis, invalidation level, next catalyst, and urgency
  \end{itemize}
\end{enumerate}

Recommended alert payload:

\begin{Shaded}
\begin{Highlighting}[]
\FunctionTok{\{}
  \DataTypeTok{"signal\_id"}\FunctionTok{:} \StringTok{"sig\_20260330\_001"}\FunctionTok{,}
  \DataTypeTok{"severity"}\FunctionTok{:} \StringTok{"ACTIONABLE"}\FunctionTok{,}
  \DataTypeTok{"event\_title"}\FunctionTok{:} \StringTok{"Will X happen by date Y?"}\FunctionTok{,}
  \DataTypeTok{"market\_id"}\FunctionTok{:} \StringTok{"12345"}\FunctionTok{,}
  \DataTypeTok{"direction"}\FunctionTok{:} \StringTok{"YES"}\FunctionTok{,}
  \DataTypeTok{"probability\_change\_15m"}\FunctionTok{:} \FloatTok{0.11}\FunctionTok{,}
  \DataTypeTok{"current\_probability"}\FunctionTok{:} \FloatTok{0.67}\FunctionTok{,}
  \DataTypeTok{"fresh\_wallet\_count\_15m"}\FunctionTok{:} \DecValTok{4}\FunctionTok{,}
  \DataTypeTok{"fresh\_wallet\_notional\_15m"}\FunctionTok{:} \DecValTok{48250}\FunctionTok{,}
  \DataTypeTok{"largest\_first\_trade\_usd"}\FunctionTok{:} \DecValTok{18000}\FunctionTok{,}
  \DataTypeTok{"cluster\_concentration"}\FunctionTok{:} \FloatTok{0.71}\FunctionTok{,}
  \DataTypeTok{"external\_public\_explanation\_score"}\FunctionTok{:} \FloatTok{0.18}\FunctionTok{,}
  \DataTypeTok{"supporting\_evidence"}\FunctionTok{:} \OtherTok{[}
    \StringTok{"4 fresh wallets entered within 11 minutes"}\OtherTok{,}
    \StringTok{"92\% of fresh{-}wallet notional is same{-}direction"}\OtherTok{,}
    \StringTok{"No comparable X/news spike yet"}
  \OtherTok{]}\FunctionTok{,}
  \DataTypeTok{"invalidation"}\FunctionTok{:} \StringTok{"probability reverts below 0.59 without fresh{-}wallet follow{-}through"}\FunctionTok{,}
  \DataTypeTok{"next\_check\_eta\_minutes"}\FunctionTok{:} \DecValTok{10}
\FunctionTok{\}}
\end{Highlighting}
\end{Shaded}

\subsection{7. ML Model Requirements}\label{7-ml-model-requirements}

\subsubsection{7.1 Modeling philosophy}\label{71-modeling-philosophy}

Use a layered modeling stack:

\begin{itemize}
\tightlist
\item
  heuristics for candidate generation,
\item
  supervised ranking for alert priority,
\item
  unsupervised anomaly detection for rare behaviors,
\item
  sequence and graph models once the feature store is mature.
\end{itemize}

\subsubsection{7.2 Required model
classes}\label{72-required-model-classes}

\begin{enumerate}
\def\labelenumi{\arabic{enumi}.}
\item
  Baseline heuristic detector

  \begin{itemize}
  \tightlist
  \item
    Threshold rules for fresh-wallet flow and probability acceleration.
  \end{itemize}
\item
  Supervised alert-ranking model

  \begin{itemize}
  \tightlist
  \item
    LightGBM or XGBoost over engineered features.
  \end{itemize}
\item
  Unsupervised anomaly detector

  \begin{itemize}
  \tightlist
  \item
    Isolation Forest, robust z-score stack, EVT tails, or one-class
    methods.
  \end{itemize}
\item
  Temporal sequence model

  \begin{itemize}
  \tightlist
  \item
    Temporal CNN, transformer, or Hawkes/state-space model for event
    episodes.
  \end{itemize}
\item
  Graph clustering model

  \begin{itemize}
  \tightlist
  \item
    Wallet-market bipartite graph for coordinated-flow communities.
  \end{itemize}
\end{enumerate}

\subsubsection{7.3 Feature families}\label{73-feature-families}

Microstructure features:

\begin{itemize}
\tightlist
\item
  price velocity
\item
  price acceleration
\item
  spread
\item
  depth imbalance
\item
  quote replenishment
\item
  large-trade share
\item
  buy/sell imbalance
\end{itemize}

Wallet features:

\begin{itemize}
\tightlist
\item
  first appearance age
\item
  first-trade notional
\item
  total traded markets
\item
  profile metadata completeness
\item
  repeat participation across related events
\item
  realized and unrealized PnL persistence if available
\end{itemize}

Context features:

\begin{itemize}
\tightlist
\item
  category
\item
  time to resolution
\item
  event popularity
\item
  public-discussion intensity
\item
  cross-market linkage strength
\end{itemize}

External evidence features:

\begin{itemize}
\tightlist
\item
  X hit count
\item
  top-post engagement
\item
  web/news hit count
\item
  corroboration lag
\item
  contradiction score
\end{itemize}

\subsubsection{7.4 Labels}\label{74-labels}

Primary labels:

\begin{itemize}
\tightlist
\item
  next-5m return
\item
  next-30m return
\item
  next-2h return
\item
  next-24h return
\end{itemize}

Research labels:

\begin{itemize}
\tightlist
\item
  time to first public corroboration
\item
  time to cross-market repricing
\item
  eventual event outcome
\item
  realized PnL from rule-based execution
\end{itemize}

Operational labels:

\begin{itemize}
\tightlist
\item
  analyst accepted / ignored / rejected
\item
  alert false positive / true positive
\end{itemize}

\subsubsection{7.5 Evaluation metrics}\label{75-evaluation-metrics}

Statistical metrics:

\begin{itemize}
\tightlist
\item
  AUROC
\item
  AUPRC
\item
  precision at k
\item
  recall at fixed false-positive rate
\item
  calibration error
\end{itemize}

Economic metrics:

\begin{itemize}
\tightlist
\item
  paper-trade PnL
\item
  Sharpe-like risk-adjusted return
\item
  hit rate net of slippage
\item
  average edge decay after alert
\end{itemize}

Operational metrics:

\begin{itemize}
\tightlist
\item
  median lead time to public corroboration
\item
  alert volume per day
\item
  false-alarm rate
\item
  time from candidate to alert
\end{itemize}

\subsubsection{7.6 Model benchmarks}\label{76-model-benchmarks}

Required baselines:

\begin{itemize}
\tightlist
\item
  raw Polymarket momentum only
\item
  order imbalance only
\item
  fresh-wallet heuristic only
\item
  external-evidence only
\item
  combined wallet-unaware microstructure model
\end{itemize}

The wallet-aware model must beat wallet-unaware baselines to justify the
project\textquotesingle s novelty claim.

\subsection{8. Backtesting and Validation
Framework}\label{8-backtesting-and-validation-framework}

\subsubsection{8.1 Point-in-time
reconstruction}\label{81-point-in-time-reconstruction}

The backtest environment must reconstruct the exact state available at
decision time.

Required:

\begin{itemize}
\tightlist
\item
  raw event stream archive
\item
  feature materialization by timestamp
\item
  event-resolution timeline
\item
  external evidence timestamps
\item
  no forward-filled profile or outcome data
\end{itemize}

\subsubsection{8.2 Validation designs}\label{82-validation-designs}

\begin{enumerate}
\def\labelenumi{\arabic{enumi}.}
\item
  Time-split validation

  \begin{itemize}
  \tightlist
  \item
    train on earlier windows, test on later windows
  \end{itemize}
\item
  Category holdout

  \begin{itemize}
  \tightlist
  \item
    test on unseen categories such as geopolitics after training on
    politics/sports
  \end{itemize}
\item
  Event-family holdout

  \begin{itemize}
  \tightlist
  \item
    keep related contracts together to avoid leakage
  \end{itemize}
\item
  Goodhart validation

  \begin{itemize}
  \tightlist
  \item
    compare performance before and after signal publication or alert
    deployment
  \end{itemize}
\end{enumerate}

\subsubsection{8.3 Event-study
framework}\label{83-event-study-framework}

For each signal episode, estimate:

\begin{itemize}
\tightlist
\item
  abnormal repricing after signal,
\item
  market-depth response,
\item
  public-corroboration lag,
\item
  cross-market propagation,
\item
  and eventual event realization.
\end{itemize}

Suggested windows:

\begin{itemize}
\tightlist
\item
  \texttt{{[}-60m,\ +5m{]}}
\item
  \texttt{{[}-60m,\ +30m{]}}
\item
  \texttt{{[}-6h,\ +24h{]}}
\item
  category-specific windows for macro and geopolitical events
\end{itemize}

\subsubsection{8.4 Paper-trading
framework}\label{84-paper-trading-framework}

Need a conservative execution simulator:

\begin{itemize}
\tightlist
\item
  use best ask for buys and best bid for sells,
\item
  include slippage based on depth,
\item
  cap size by market liquidity,
\item
  stop new entries near resolution if liquidity is weak,
\item
  and separate model score from execution assumptions.
\end{itemize}

\subsubsection{8.5 Success criteria}\label{85-success-criteria}

Minimum thesis-quality success standard:

\begin{itemize}
\tightlist
\item
  statistically significant improvement over wallet-unaware baselines,
\item
  positive paper-trading edge after slippage,
\item
  and meaningful lead time over public corroboration on a nontrivial
  subset of alerts.
\end{itemize}

\subsection{9. Risks, Limitations, and Ethical
Considerations}\label{9-risks-limitations-and-ethical-considerations}

\subsubsection{9.1 Identification risk}\label{91-identification-risk}

Fresh wallets are not fresh humans.

You should never label a wallet cluster as "insider" in product outputs.
Use:

\begin{itemize}
\tightlist
\item
  anomalous,
\item
  coordinated,
\item
  fresh-wallet,
\item
  potential informed flow.
\end{itemize}

\subsubsection{9.2 False accusation
risk}\label{92-false-accusation-risk}

Publicly naming wallets or implying illegal conduct creates legal and
ethical risk.

Mitigation:

\begin{itemize}
\tightlist
\item
  keep wallet identifiers internal,
\item
  hash identifiers in user-facing alerts by default,
\item
  and separate anomaly detection from any legal conclusion.
\end{itemize}

\subsubsection{9.3 Market-impact risk}\label{93-market-impact-risk}

If your own alerting becomes widely acted upon, it may move prices and
partially destroy the signal you are studying.

This is not just a trading issue; it is part of H1.

\subsubsection{9.4 Data and platform
risk}\label{94-data-and-platform-risk}

\begin{itemize}
\tightlist
\item
  API contracts may change.
\item
  Rate limits may tighten.
\item
  Websocket event formats may drift.
\item
  Public profile metadata may be incomplete or noisy.
\end{itemize}

\subsubsection{9.5 Regulatory and compliance
risk}\label{95-regulatory-and-compliance-risk}

Prediction-market execution legality varies by jurisdiction and platform
policy.

This SRS should treat automated execution as a separate optional layer.
The core system should be a research, monitoring, and alerting engine
first.

\subsubsection{9.6 Statistical risk}\label{96-statistical-risk}

\begin{itemize}
\tightlist
\item
  rare-event overfitting,
\item
  selection bias,
\item
  lookahead leakage,
\item
  and survivorship bias are major threats.
\end{itemize}

\subsubsection{9.7 Ethical framing}\label{97-ethical-framing}

The project should be framed as:

\begin{itemize}
\tightlist
\item
  market intelligence,
\item
  anomaly detection,
\item
  and information-aggregation research,
\end{itemize}

not as an accusation engine or surveillance system targeting
individuals.

\subsection{10. Development Roadmap}\label{10-development-roadmap}

\subsubsection{Phase 0: Stabilize the
collector}\label{phase-0-stabilize-the-collector}

Target: 1-2 weeks

\begin{itemize}
\tightlist
\item
  Remove or parameterize sports-only filtering.
\item
  Fix trades endpoint usage to use documented condition IDs.
\item
  Add both-asset coverage where needed.
\item
  Add ingestion health metrics.
\item
  Resume continuous collection in a stable environment.
\end{itemize}

Deliverable:

\begin{itemize}
\tightlist
\item
  reliable fresh data across target categories.
\end{itemize}

\subsubsection{Phase 1: Add wallet-aware
ingestion}\label{phase-1-add-wallet-aware-ingestion}

Target: 2-3 weeks

\begin{itemize}
\tightlist
\item
  extend schema for wallets, activity, and positions
\item
  store \texttt{proxyWallet}, \texttt{conditionId},
  \texttt{transactionHash}, \texttt{outcome}
\item
  add profile enrichment
\item
  archive raw payloads
\end{itemize}

Deliverable:

\begin{itemize}
\tightlist
\item
  account-aware market dataset suitable for H2.
\end{itemize}

\subsubsection{Phase 2: Build the signal
engine}\label{phase-2-build-the-signal-engine}

Target: 3-4 weeks

\begin{itemize}
\tightlist
\item
  implement fresh-wallet detector
\item
  implement coordinated-flow episode builder
\item
  add microstructure anomaly features
\item
  create first ranked signal score
\end{itemize}

Deliverable:

\begin{itemize}
\tightlist
\item
  offline signal table with replayable episodes.
\end{itemize}

\subsubsection{Phase 3: Build external evidence and
alerting}\label{phase-3-build-external-evidence-and-alerting}

Target: 2-4 weeks

\begin{itemize}
\tightlist
\item
  integrate X recent search
\item
  integrate web/news search
\item
  add Telegram alerts
\item
  add suppression and severity logic
\end{itemize}

Deliverable:

\begin{itemize}
\tightlist
\item
  live alert loop with evidence summaries.
\end{itemize}

\subsubsection{Phase 4: Backtesting and paper
trading}\label{phase-4-backtesting-and-paper-trading}

Target: 4-6 weeks

\begin{itemize}
\tightlist
\item
  point-in-time replay
\item
  event studies
\item
  baseline comparisons
\item
  conservative execution simulator
\end{itemize}

Deliverable:

\begin{itemize}
\tightlist
\item
  first empirical answer to H1/H2/H3.
\end{itemize}

\subsubsection{Phase 5: Model refinement and thesis
packaging}\label{phase-5-model-refinement-and-thesis-packaging}

Target: 4-6 weeks

\begin{itemize}
\tightlist
\item
  supervised ranking model
\item
  graph clustering for coordinated wallets
\item
  ablations and robustness checks
\item
  final research memo and visuals
\end{itemize}

Deliverable:

\begin{itemize}
\tightlist
\item
  thesis-grade results, figures, and defensible methodology.
\end{itemize}

\subsection{Suggested Immediate Engineering
Priorities}\label{suggested-immediate-engineering-priorities}

\begin{enumerate}
\def\labelenumi{\arabic{enumi}.}
\tightlist
\item
  Fix the collector so it covers non-sports markets and stores
  documented trade fields.
\item
  Move from SQLite-only local storage to a cloud Postgres plus raw
  archive design.
\item
  Add wallet-aware tables before writing any serious anomaly model.
\item
  Preserve enough closed-market microstructure to run event studies.
\item
  Build the alerting and external-evidence layer only after the raw
  market data is trustworthy.
\end{enumerate}

\subsection{Sources}\label{sources}

Primary platform documentation:

\begin{itemize}
\tightlist
\item
  Polymarket trades endpoint:
  \url{https://docs.polymarket.com/api-reference/core/get-trades-for-a-user-or-markets}
\item
  Polymarket user activity:
  \url{https://docs.polymarket.com/api-reference/core/get-user-activity}
\item
  Polymarket current positions for user:
  \url{https://docs.polymarket.com/api-reference/core/get-current-positions-for-a-user}
\item
  Polymarket positions for a market:
  \url{https://docs.polymarket.com/api-reference/core/get-positions-for-a-market}
\item
  Polymarket public profile:
  \url{https://docs.polymarket.com/api-reference/profiles/get-public-profile-by-wallet-address}
\item
  Polymarket total markets traded:
  \url{https://docs.polymarket.com/api-reference/misc/get-total-markets-a-user-has-traded}
\item
  Polymarket trader leaderboard:
  \url{https://docs.polymarket.com/api-reference/core/get-trader-leaderboard-rankings}
\item
  Polymarket market websocket channel:
  \url{https://docs.polymarket.com/market-data/websocket/market-channel}
\end{itemize}

Cloud platform documentation:

\begin{itemize}
\tightlist
\item
  AWS Lambda with SQS:
  \url{https://docs.aws.amazon.com/lambda/latest/dg/with-sqs.html}
\item
  AWS SQS error handling for Lambda:
  \url{https://docs.aws.amazon.com/lambda/latest/dg/services-sqs-errorhandling.html}
\item
  Cloud Run WebSockets:
  \url{https://cloud.google.com/run/docs/triggering/websockets}
\item
  Cloud Run minimum instances:
  \url{https://cloud.google.com/run/docs/configuring/min-instances}
\item
  Cloud Run request timeout:
  \url{https://cloud.google.com/run/docs/configuring/request-timeout}
\item
  Cloud Run scheduling with Cloud Scheduler:
  \url{https://cloud.google.com/run/docs/triggering/using-scheduler}
\item
  Railway cron jobs:
  \url{https://docs.railway.com/reference/cron-jobs?hidden=true}
\item
  Railway data and storage: \url{https://docs.railway.com/data-storage}
\item
  Railway volume reference:
  \url{https://docs.railway.com/reference/volumes}
\end{itemize}

Alerting and search documentation:

\begin{itemize}
\tightlist
\item
  Telegram Bot API: \url{https://core.telegram.org/bots/api}
\item
  Discord webhooks:
  \url{https://discord.com/developers/docs/resources/webhook}
\item
  ntfy publishing: \url{https://docs.ntfy.sh/publish/}
\item
  Twilio messages API:
  \url{https://www.twilio.com/docs/messaging/api/message-resource}
\item
  Google Custom Search JSON API overview:
  \url{https://developers.google.com/custom-search/v1/overview}
\item
  X recent search API:
  \url{https://docs.x.com/x-api/posts/search-recent-posts}
\end{itemize}

Research references:

\begin{itemize}
\tightlist
\item
  Wolfers, Zitzewitz, "Prediction Markets" (NBER 2004):
  \url{https://www.nber.org/papers/w10504}
\item
  Wolfers, Zitzewitz, "Prediction Markets in Theory and Practice" (NBER
  2006): \url{https://www.nber.org/papers/w12083}
\item
  Snowberg, Wolfers, Zitzewitz, "How Prediction Markets Can Save Event
  Studies" (NBER 2011): \url{https://www.nber.org/papers/w16949}
\item
  Snowberg, Wolfers, Zitzewitz, "Prediction Markets for Economic
  Forecasting" (NBER 2012): \url{https://www.nber.org/papers/w18222}
\item
  Oprea, Hanson, Porter, "Affecting Policy by Manipulating Prediction
  Markets: Experimental Evidence":
  \url{https://www.sciencedirect.com/science/article/pii/S0167268112002223}
\item
  Yang, Zhu, "Back-Running: Seeking and Hiding Fundamental Information
  in Order Flows":
  \url{https://papers.ssrn.com/sol3/papers.cfm?abstract_id=2583915}
\item
  Easley, Lopez de Prado, O\textquotesingle Hara, "The Volume Clock:
  Insights into the High Frequency Paradigm":
  \url{https://papers.ssrn.com/sol3/papers.cfm?abstract_id=2034858}
\item
  Reichenbach, Walther, "Exploring Decentralized Prediction Markets:
  Accuracy, Skill, and Bias on Polymarket":
  \url{https://papers.ssrn.com/sol3/papers.cfm?abstract_id=5910522}
\item
  Ng, Peng, Tao, Zhou, "Price Discovery and Trading in Modern Prediction
  Markets":
  \url{https://papers.ssrn.com/sol3/papers.cfm?abstract_id=5331995}
\item
  Burgi, Deng, Whelan, "Makers and Takers: The Economics of the Kalshi
  Prediction Market":
  \url{https://papers.ssrn.com/sol3/papers.cfm?abstract_id=5502658}
\end{itemize}

\end{document}
